Die innere Energie eines idealen Gases ist mit der Zustandsgleichung $p\,V = nRT$
\begin{align}
 U &= \frac{f}{2}\,nRT = \frac{f}{2}\,p\,V
\end{align}
Sie ist also proportional zum Produkt $p\,V$. Neu ist $V' = \left(1+\fVO\right)V$ und $p' = \left(1-\fpO\right)p$, die relative Änderung ist
\begin{align}
 \frac{\Delta U}{U} &= \frac{p'V'}{p\,V}-1 = \dUrF \\
   &= \left(1+\fVO\right)\cdot\left(1-\fpO\right)-1 \\
   &= \dUr \approx \result{\dUrP{0}{2}}
\end{align}
Die innere Energie nimmt ab, und mit ihr die Temperatur.
